\section{Canonical Structure of the Latent Cognition Research Program}

The Latent Cognition Trilogy should be read as the main theoretical sequence within a broader research program on latent cognition, retrieval-native intelligence, and bounded artificial agency. The Trilogy is not an isolated manuscript collection, nor is it a container for every related hypothesis, architectural proposal, or companion paper. Rather, it functions as the central narrative and conceptual axis around which a wider body of work is organized.

Within this structure, \textit{Differentiable Latent Indexing in LLMs: Toward Native Retrieval as a Cognitive Function} serves as the architectural foundation of the research program. That work develops the systems-level premise that retrieval should not remain an external augmentation layer, as in conventional retrieval-augmented generation architectures, but may instead be modeled as a differentiable cognitive function embedded within the representational dynamics of transformer-based systems. For this reason, the DLI-LLM paper should be understood as the architectural substrate beneath the Trilogy rather than as a supplementary appendix.

The three companion hypothesis papers occupy a different position. They do not define the base architecture of the program. Instead, they formalize adjacent theoretical claims that support, constrain, or extend the main conceptual sequence developed in the Trilogy. Their purpose is to preserve analytical precision around claims that would otherwise overload the central manuscript.

Accordingly, the present work adopts the following canonical structure: the DLI-LLM paper functions as the foundational architecture paper; the Latent Cognition Trilogy forms the main theoretical sequence; and the companion hypothesis papers operate as formal satellite works within the same research program. This organization allows the Trilogy to remain coherent while preserving the broader intellectual architecture required for future papers, extensions, and domain-specific applications.
